\chapter{The physical small-$1-r$ source has a forced scale bound}
\label{component:p0c68f3aecc59}
\noindent\textbf{Component }\path{p0c68f3aecc59}\par
{\small\noindent Source: \path{R13/advances/weak_small_gap_scale/boundary_scale_lemma.tex}\par}
{\small\noindent SHA256: \path{0c68f3aecc59fba087d199c915d7c464be358b028e41022ad9704248c96e6860}\par}
{\small\noindent Inventory use: BS1-BS5\par}
\medskip
\section{The physical small-$1-r$ source has a forced scale bound}
\label{p0c68f3aecc59:sec:r12-small-gap-scale}

Retain the original physical C2 source, its exact contact and strict
gaps, with $n\ge0$, $d=1$, $B=H+L$, and $B,z,y>0$
infinitesimal. Put
\[
 a=z/B,\quad t=y/B,\quad h=H/B,\quad e=1-r>0.
\]
Assume $e,h$ infinitesimal and $a,t,G/B,j,L/B$ positive bounded
units. Thus $\rho=\operatorname{st}(a)>0$,
$\tau=\operatorname{st}(t)>1$, and $r$ is a unit.

\begin{lemma}
The exact original source identity is
\[
 2r\Delta n=eQ-BK,\qquad \Delta=\delta/B,             \tag{BS1}
\]
where
\[
\begin{split}
 K&=(1+a)(h+t)(1+aB),\\
 Q&=1+B[1+t+ah+2ar+at]\\
  &\quad+B^2\{a(h+t)+a^2[h(r+1)+t(1+r-r^2)]\}.
\end{split}                                            \tag{BS2}
\]
Both $Q,K$ are positive bounded units, with residues
$Q_0=1$, $K_0=(1+\rho)\tau$, and $\Delta_0=1+\rho$.
Consequently source nonnegativity is exactly equivalent to
\[
 0<B/e\le Q/K,\qquad
 0\le n/e\le Q/(2r\Delta).                            \tag{BS3}
\]
In particular $B/e$ and $n/e$ are bounded, regardless of the
relative rate of $h$.
\end{lemma}
\begin{proof}
The original normalized source formula is
\[
 2r\Delta n=(1+ar)(1+aBr)[1+B(1-h+ar)]-r\Delta,
\]
with
\[
\begin{split}
 \Delta={}&a+(1+aB)+(1+aB)B+2aB(1+aB)r+a^3B^2r^2
       +2a^2Br\\
       &+tB[1+a+aB+a^2Br(2-r)].
\end{split}
\]
Subtracting its value at $r=1$ and dividing by $1-r$ gives
the polynomial $Q$ in (BS2); the value at $r=1$ is $-BK$.
Every term in $Q$ is positive, since $0<r<1$ implies
$1+r-r^2>0$. The stated residues and inequalities follow.
\end{proof}

\paragraph{The literal $r=1$ face is not an admissible anchor.}
At $r=1$ and still $B>0$, the exact identity becomes
\[
 n\big|_{r=1}
 =-\frac{B(h+t)(1+aB)}
 {2[1+B(1+2a+t)+aB^2(1+a+t)]}<0.                      \tag{BS4}
\]
Thus simply setting $r=1$ exits the original nonnegative-source
domain. An event-face theorem requiring $n\ge0$ cannot be applied
to that unmodified endpoint. This is a limitation of that endpoint
reduction, not a refutation of any alternative full-parent proof.

\paragraph{What this repairs, and what remains unresolved.}
Let $\lambda_*=B/e$, which is bounded by (BS3). Its residue may
be zero. Joint substitution $B=e\lambda_*$, $r=1-e$ gives
\[
 (n/e)_0=\frac{1-\lambda_{*,0}(1+\rho)\tau}{2(1+\rho)},
 \qquad (m/e)_0=\frac{\rho^2}{2(1+\rho)}>0.           \tag{BS5}
\]
For the scalar coefficient $\beta_s$ in (WA4),
\[
 (\beta_s/e)_0
 =\rho^2\left[\frac1{2(1+\rho)}
                   +\frac{\lambda_{*,0}\tau}{2(\tau-1)}\right]>0.
                                                               \tag{BS6}
\]
So source physicality removes the apparent diagonal-coefficient
pole after scaling $R_2$ by $\sqrt e$.
It does not remove the previously established cross-coefficient
poles: the same exact source still forces
\[
 (j_0)_0=\frac{\rho\tau(\tau+1)}{4(\tau-1)},\quad
 (j_1)_0=\frac{\rho\tau^2}{4(\tau-1)},\quad
 (f_s)_0=-\frac{\rho\tau}{4}.
\]
Thus that fixed Gram still has unbounded $J_s/\sqrt e$ and
$f_s/\sqrt e$. No new local full-parent certificate or all-tests
admission is claimed here. The decisive result is the exact
physical scale bound and the exclusion of the literal $r=1$
nonnegative-source anchor. Source physicality and contact are
retained as hypotheses throughout; no coefficient-sign argument
is used to infer global physicality.
